\langsection[subsection]{Know Your Responsibilities}{Знать свою зону ответственности}

\begin{langblock}{0.425}
	\begin{eng}
	``Accident`` implies an unforeseen event that occurs without anyone's fault or negligence. Most often in traffic, that is not the case. In fact, most people involved in a crash can usually claim some responsibility for what took place.
	
	Consider a situation where someone tries to squeeze through an intersection on a yellow light turning red. Your light turns green. You pull into the intersection without checking for possible latecomers. That is all it takes for the two of you to tangle. It was the other driver's responsibility to stop. And it was your responsibility to look before pulling out. Neither of you held up your end of the deal. Just because someone else is the first to start the chain of events leading to a crash, doesn't leave any of us free of responsibility.
	
	As a rider, you can't be sure that other operators will see you or yield the right of way. To lessen your chances of a crash occurring:
	\begin{itemize}
		\item \textbf{Be visible} --- wear proper clothing, use your headlight, ride in the best lane position to see and be seen.
		\item \textbf{Communicate your intentions} --- use the proper signals, brake light, and lane position.
		\item \textbf{Maintain an adequate space cushion} --- following, being followed, lane sharing, passing and being passed.
		\item \textbf{Search your path} of travel 12 seconds ahead.
		\item \textbf{Identify and separate} multiple hazards.
		\item \textbf{Be prepared to act} --- remain alert and know how to carry out proper crash-avoidance skills.
	\end{itemize}
	Blame doesn't matter when someone is injured in a crash. There is rarely a single cause of any crash. The ability to ride aware, make critical decisions, and carry them out separates responsible riders from all the rest. Remember, it is up to you to keep from being the cause of, or an unprepared participant in, any crash.
	\end{eng}
	\begin{rus}
		Термин «авария» подразумевает непредвиденное событие, произошедшее без чьей-либо вины или халатности. Но чаще всего в дорожном движении это не так. На самом деле большинство участников столкновения, как правило, несут определенную долю ответственности за произошедшее. 
		
		Рассмотрим ситуацию, когда кто-то пытается проскочить перекресток на желтый сигнал, сменяющийся красным. Ваш сигнал переключается на зеленый. Вы выезжаете на перекресток, не проверив наличие тех, кто пытался проскочить. Этого достаточно, чтобы произошло столкновение. Обязанность остановиться лежала на другом водителе. Но вашей обязанностью было посмотреть по сторонам, прежде чем выезжать. Ни один из вас не выполнил свои обязательства. Тот факт, что кто-то первым положил начало цепочке событий, приведших к столкновению, не освобождает никого из вас от ответственности.
		
		Как водитель мотоцикла, вы не можете быть уверены в том, что другие участники движения заметят вас или уступят дорогу. Чтобы снизить вероятность возникновения ДТП:
		\begin{itemize}
			\item \textbf{Будьте видимы} "--- используйте надлежащую одежду, включайте фару, занимайте оптимальное положение в полосе для обзора и заметности.
			\item \textbf{Сообщайте о своих намерениях} "--- используйте соответствующие сигналы, стоп-сигнал и положение в полосе.
			\item \textbf{Поддерживайте достаточную дистанцию} "--- при движении за кем-либо, когда кто-то следует за вами, при совместном использовании полосы, при обгоне и когда обгоняют вас.
			\item \textbf{Контролируйте} пространство на 12 секунд вперед.
			\item \textbf{Выявляйте и разделяйте} множественные источники опасности.
			\item \textbf{Будьте готовы к действию} "--- сохраняйте бдительность и владейте навыками предотвращения столкновений.
		\end{itemize}
		Упреки не имеет значения, когда в ДТП уже кто-то получает травмы. Редко какое-либо столкновение имеет единственную причину. Способность осознанно управлять, принимать и выполнять обоснованные решения отличает ответственных водителей от всех остальных. Помните: именно от вас зависит, не станете ли вы причиной ДТП и будете ли готовы его предотвратить.
\end{rus}
\end{langblock}
